\chapter{General Relativity and Gravitational Waves}
\label{ch:general_relativity}


\begin{objectivebox}
\begin{itemize}
    \item Re-derive the linearised Einstein field equations as wave equations on a trace-reversed LC transmission line.
    \item Understand spacetime curvature as a continuous impedance gradient in the $\mathcal{M}_A$ dielectric network.
    \item Compute gravitational time dilation and redshift from the local refractive index $n(r) = 1/S(r)$.
\end{itemize}
\end{objectivebox}

\section{The Ontology of Spacetime Curvature}
Einstein's General Relativity (GR) is a comprehensive framework in differential geometry. It models gravity not as a direct force, but as the curvature of a 4-dimensional Spacetime manifold caused by the presence of mass and energy. 

A long-standing interpretive question in GR is the physical ontology of the manifold itself. The metric carries momentum, possesses effective inertia, and transmits waves at finite speed ($c$)---properties typically associated with a physical medium rather than pure geometry.

Applied Vacuum Engineering (AVE) resolves this ontological paradox by defining ``Curved Spacetime'' as the variable scalar Capacitance ($C$) and Inductance ($L$) of a structured, continuous dielectric super-fluid.

\section{The Stress-Energy Tensor as LC Energy Density}
The core of General Relativity is Einstein's Field Equation:
\begin{resultbox}{Einstein Field Equation}
\begin{equation}
    R_{\mu\nu} - \frac{1}{2}R g_{\mu\nu} + \Lambda g_{\mu\nu} = \frac{8\pi G}{c^4} T_{\mu\nu}
\end{equation}
\end{resultbox}

In the AVE framework, the Stress-Energy Tensor ($T_{\mu\nu}$) on the right side of the equation is the classical Electromagnetic Energy Density ($U$) of the local LC vacuum:
\begin{resultbox}{Stress-Energy as LC Energy Density}
\begin{equation}
    T_{\mu\nu} \equiv U_{\mu\nu} = \frac{1}{2}\epsilon_0 |\mathbf{E}|^2 + \frac{1}{2}\mu_0 |\mathbf{H}|^2
\end{equation}
\end{resultbox}

Furthermore, the mathematical Metric Tensor ($g_{\mu\nu}$) describing the curvature on the left side of the equation is isomorphic to the macroscopic structural variable impedance parameters ($\varepsilon_{eff}, \mu_{eff}$) of the dielectric matrix. Under \textbf{Symmetric Gravity}, both constitutive parameters scale with the refractive index $n(r) = 1 + 2GM/(rc^2)$:
\begin{resultbox}{Symmetric Gravity Impedance}
\begin{equation}
    \varepsilon_{eff}(r) = \varepsilon_0 \cdot n(r), \qquad
    \mu_{eff}(r) = \mu_0 \cdot n(r), \qquad
    Z(r) = \sqrt{\frac{\mu_{eff}}{\varepsilon_{eff}}} = Z_0 \;\text{(invariant)}
\end{equation}
\end{resultbox}
The Schwarzschild radius $r_s = 2GM/c^2$ marks the point where the local strain reaches the Axiom~4 saturation limit ($S \to 0$): both $\mu_{eff}$ and $\varepsilon_{eff}$ collapse to zero, the local impedance drops to zero ($Z \to 0$), and the reflection coefficient reaches $\Gamma = -1$ (total confinement). This is the physical Event Horizon---not a coordinate singularity, but a \textbf{catastrophic impedance boundary}.

When localised topological energy (mass) is present, it draws continuous phase-locked energy from the surrounding LC grid. This creates an inductive deficit in the adjacent vacuum, analogous to a density gradient in fluid dynamics. This impedance gradient ($Z = \sqrt{\mu/\epsilon}$) acts as an optical refractive index, bending the propagation trajectories of passing light and accelerating other mass-bearing geometric knots down the gradient. ``Gravity'' is macroscopic dielectric refraction.

\section{Gravitational Waves as Inductive Shear}
In 2015, LIGO detected "Gravitational Waves" from merging black holes. Mainstream physics describes this as "ripples in the fabric of spacetime itself."

In the AVE framework, a black hole corresponds to a localized region of maximum dielectric saturation where the LC grid reaches its capacitive yield point (the Event Horizon). When two such massive topological stress-concentrations orbit each other, they act as macroscopic impellers driving transverse shear waves through the electro-mechanical condensate.

\begin{figure}[H]
    \centering
    \includegraphics[width=0.85\textwidth]{gravitational_waves_lc_static.pdf}
    \caption{A 2D FDTD simulation of two super-massive topological nodes orbiting in a binary pair. Their rotational acceleration acts as an impeller, dragging the local LC grid. This mechanical pumping action radiates transverse displacement current ($d\vec{D}/dt$) shear-waves outward into the cosmos. These are ``Gravitational Waves''---identical in mathematical description to standard acoustic shear-waves propagating through an elastic crystalline matrix.}
    \label{fig:gravitational_waves}
\end{figure}

In this interpretation, gravitational waves are low-frequency macroscopic inductive strain-waves propagating through the structured LC condensate.

By identifying the vacuum as a physical, variable-impedance LC medium, General Relativity is placed in direct correspondence with classical Continuum Mechanics and Electrodynamics. The unification of gravity with quantum field theory reduces to recognising the electromagnetic character of the spacetime substrate.

\subsection*{Gravitational Wave Regime Classification}
\begin{center}
\small
\renewcommand{\arraystretch}{1.2}
\begin{tabular}{lccl}
\toprule
\textbf{Source} & \textbf{Strain $h$} & \textbf{$V_{GW}/V_{snap}$} & \textbf{Regime} \\
\midrule
GW150914 (BBH) & $10^{-21}$ & $\sim 10^{-28}$ & I (deeply linear) \\
GW170817 (BNS) & $10^{-22}$ & $\sim 10^{-29}$ & I (deeply linear) \\
Pulsar timing  & $10^{-15}$ & $\sim 10^{-22}$ & I (linear) \\
Near-merger ($r \lesssim 10\,r_s$) & $10^{-1}$ & $\sim 10^{-8}$ & I--II boundary \\
\bottomrule
\end{tabular}
\end{center}

\noindent All observed GW signals operate far below the dielectric yield threshold. The vacuum responds linearly, producing lossless propagation at speed $c$---the mathematical predictions of linearized GR emerge exactly. The nonlinear regime (II--III) is only reached in the near-merger zone where waveform corrections become significant.

\section{The General Relativistic Lense-Thirring Manifold}
The ultimate test of any unified acoustic theory of gravity is predicting the extreme geometric deformation surrounding a rapidly rotating supermassive black hole. In mainstream physics, this requires solving the complex tensor geodesics of the \textit{Kerr Metric}. 

However, within the AVE framework, curved spacetime geometry does not exist. The physics engine simulation of the \textit{Interstellar Gargantua Black Hole} (Mass $\sim 10^8 M_\odot$, Spin $a \approx 0.999$) relies strictly on \textbf{Op14 Topological Saturation}. 

The abstract GR Riemannian manifold $g_{\mu\nu}$ collapses directly into a macroscopic impedance manifold:
\begin{enumerate}
    \item \textbf{A Macroscopic Index of Refraction ($n$):} The presence of mass creates a spherical density gradient in the LC matrix. In isotropic coordinates, the effective refractive index governing the speed of electromagnetic waves is:
    \begin{resultbox}{Gravitational Refractive Index Gradient}
\begin{equation}
        n(r) = \frac{\left(1 + \frac{r_s}{2r}\right)^3}{1 - \frac{r_s}{2r}}
\end{equation}
\end{resultbox}
    where $r_s$ is the Schwarzschild radius equivalent. As $r \to r_s/2$, the matrix yield tension approaches infinity ($n \to \infty$), acting as an absolute classical optical trap.
    
    \item \textbf{Topological Frame-Dragging (Lense-Thirring):} The rotation of the mass singularity creates an asymmetric strain footprint in the surrounding dielectric matrix. Unlike standard aerospace models, this is not a circulating fluid current. The angular momentum convolutes the static $Z_{local}$ impedance gradient into a non-reciprocal topological bias. In a Kerr spacetime with angular momentum parameter $a = a_\star M$, the saturation metric acquires an angular dependence:
    \begin{resultbox}{Frame-Dragging Impedance Convolution}
\begin{equation}
        \omega(r) = \frac{2Mar}{(r^2 + a^2)^2}
\end{equation}
\end{resultbox}
    which is strongest at the horizon and decays as $\sim r^{-3}$ at large radii, generating differential phase-delays for co-rotating versus counter-rotating light.
\end{enumerate}

By applying \textbf{Hamiltonian Symplectic Euler Integration} to track the momentum $\vec{p}$ of photon rays backwards through this refractive impedance gradient, the Python CFD engine derives the ``bent accretion disk'' architecture described by Kip Thorne.

The simulation renders a Shakura--Sunyaev thin accretion disk ($T_{\text{in}} = 10\,000$\,K at the inner edge, $r_{\text{in}} = 1.0\,r_{\text{iso}}$) around a near-extremal Kerr-analogue defect ($a_\star = 0.999$). Three distinct observational signatures emerge naturally from the topological model:
\begin{itemize}
    \item \textbf{The D-shaped shadow}: The spinning mass acts as a macroscopic topological impeller generating an asymmetric impedance gradient. Rays traversing the retrograde side encounter a stricter Op14 saturation profile, increasing their refractive capture radius. Conversely, rays on the prograde side propagate through a mechanically relaxed tensor, allowing them to graze closer to the horizon before capture. This differential refractive trap flattens the shadow boundary on the prograde side, producing the characteristic ``D-shape'' predicted by the Kerr metric without continuous manifolds.
    \item \textbf{Relativistic Doppler beaming}: The approaching side of the orbiting accretion disk is blue-shifted and brightened ($I_{\text{obs}} = I_{\text{emit}} \times \delta^3$ for optically thick emission).
    \item \textbf{Gravitational redshift}: The isotropic-coordinate redshift factor $z = W/U = (1 + r_h/r)/(1 - r_h/r)$ suppresses emission from the innermost disk annuli, producing the dark gap between the photon ring and the bright disk body.
\end{itemize}

\begin{figure}[H]
    \centering
    \includegraphics[width=1.0\textwidth]{gargantua_acoustic_vortex.png}
    \caption{\textbf{Hamiltonian Raymarching of the Impedance Manifold} ($a_\star = 0.999$). By mapping the Kerr metric into an asymmetric refractive gradient ($n$), numerical Hamiltonian optics reproduce the geometry of the Gargantua black hole. The shadow exhibits the characteristic D-shape from prograde strain asymmetry, while the disk brightness asymmetry demonstrates classical relativistic Doppler beaming. No non-Euclidean manifold curvature is required.}
    \label{fig:gargantua}
\end{figure}

\section{K4-TLM Gravitational Lensing Cross-Validation}
\label{sec:k4_tlm_lensing}

The preceding sections derive photon deflection analytically---$\delta = 4GM/(bc^2)$---and validate it via Hamiltonian raymarching. We now cross-validate the same physics with the K4-TLM native lattice dynamics simulator (Section~\ref{sec:k4_tlm}, Vol.~4), which time-evolves electromagnetic waves natively through the pure Diamond lattice vacuum.

\subsection{Method}

In a true K4-TLM Diamond lattice, the presence of mass is implemented strictly via Op14 continuous saturation. An exceptionally dense coordinate is designated as the mass origin. As incident scalar amplitude $\sum|V|$ grows uncontrollably near this point, Op14 dictates that the local geometry saturates. 

This builds a static $Z_{local}$ gradient across the network:
\begin{equation}
    Z_{\text{local}}(r) = \frac{Z_0}{S^{1/4}}
\end{equation}
The refractive index operates directly within the explicit 4-port scattering matrices ($S_{ij} = \frac{y_{local}}{2} - \delta_{ij}$), passively slowing the effective wave speed $c_{\text{local}} = c_0 \sqrt{Z_0/Z_{\text{local}}}$.

A photon wavepacket is injected into the lattice, propagating toward the mass defect with a finite impact parameter.

\subsection{Results}

As the wide photon wavefront encounters the topological saturation gradient, the inner side of the wavefront (closest to the mass) enters the high-impedance zone first. This forces the inner pulse to slow down, while the outer pulse remains traveling at $c_0$ in the free vacuum. 

This impedance mismatch organically pivotal the entire wavefront downward into the mass defect.

\begin{figure}[ht]
    \centering
    \includegraphics[width=1.0\textwidth]{gravity_lensing_proof.png}
    \caption{\textbf{K4-TLM Gravitational Lensing Native Output.} An explicit time-domain solve of a photon wavepacket refracting through a saturated Diamond Lattice region. The wavepacket hooks downwards, confirming that astrophysical ``Gravitational Lensing'' perfectly isomorphic to discrete optical refraction in a mathematically pure LC substrate.}
    \label{fig:k4_tlm_lensing_time}
\end{figure}

\section{GR $\leftrightarrow$ AVE Translation Dictionary}
\label{sec:gr_ave_translation}

The preceding sections demonstrate that General Relativity emerges naturally from the variable-impedance LC medium.  Table~\ref{tab:trans_gravity} makes the ontological mapping explicit.

% ═══════════════════════════════════════════════════════════════
% DOMAIN TRANSLATION TABLE: General Relativity / Gravity ↔ AVE
% Canonical location: manuscript/common/translation_gravity.tex
% ═══════════════════════════════════════════════════════════════
%
% Usage: % ═══════════════════════════════════════════════════════════════
% DOMAIN TRANSLATION TABLE: General Relativity / Gravity ↔ AVE
% Canonical location: manuscript/common/translation_gravity.tex
% ═══════════════════════════════════════════════════════════════
%
% Usage: % ═══════════════════════════════════════════════════════════════
% DOMAIN TRANSLATION TABLE: General Relativity / Gravity ↔ AVE
% Canonical location: manuscript/common/translation_gravity.tex
% ═══════════════════════════════════════════════════════════════
%
% Usage: \input{../../common/translation_gravity.tex}

\begin{table}[H]
\centering
\caption{General Relativity $\leftrightarrow$ AVE Translation.  Gravity is macroscopic dielectric refraction; no curved manifold is required.}
\label{tab:trans_gravity}
\small
\renewcommand{\arraystretch}{1.3}
\begin{tabular}{@{}p{3.8cm} p{4.0cm} p{5.8cm}@{}}
\toprule
\textbf{General Relativity} & \textbf{AVE Equivalent} & \textbf{Relationship} \\
\midrule
Metric tensor $g_{\mu\nu}$ & Variable $\varepsilon_{eff}, \mu_{eff}$ & $n(r) = 1 + 2GM/(rc^2)$; symmetric impedance gradient. \\
Spacetime curvature & Dielectric refraction & Geodesic $\equiv$ Fermat path through refractive medium. \\
Stress-energy $T_{\mu\nu}$ & LC energy density $U$ & $U = \tfrac{1}{2}\varepsilon_0 |E|^2 + \tfrac{1}{2}\mu_0 |H|^2$. \\
Schwarzschild radius $r_s$ & Saturation boundary ($S \to 0$) & $Z \to 0, \; \Gamma = -1$: total impedance collapse. \\
Event horizon & Catastrophic impedance boundary & Axiom~4: dielectric yield $\to$ total confinement. \\
Gravitational waves & Transverse inductive shear waves & Low-frequency shear modes in the LC condensate. \\
Frame dragging (Kerr) & Acoustic vortex flow $\vec{v}_\phi$ & Gordon metric: $\omega(r) = 2Mar/(r^2+a^2)^2$. \\
Gravitational lensing & Optical refraction ($n > 1$) & $n(r) = (1+r_s/2r)^3/(1-r_s/2r)$. \\
$G$ (Newton) & TT strain projection & $G = \alpha^2 c^3 \ell_{node}/(28\pi m_e)$ (Axiom~2). \\
``Dark matter'' halo & Regime~I vacuum drag & $v_{flat} = (GM a_0)^{1/4}$; $a_0 = cH_\infty/(2\pi)$. \\
\bottomrule
\end{tabular}
\end{table}


\begin{table}[H]
\centering
\caption{General Relativity $\leftrightarrow$ AVE Translation.  Gravity is macroscopic dielectric refraction; no curved manifold is required.}
\label{tab:trans_gravity}
\small
\renewcommand{\arraystretch}{1.3}
\begin{tabular}{@{}p{3.8cm} p{4.0cm} p{5.8cm}@{}}
\toprule
\textbf{General Relativity} & \textbf{AVE Equivalent} & \textbf{Relationship} \\
\midrule
Metric tensor $g_{\mu\nu}$ & Variable $\varepsilon_{eff}, \mu_{eff}$ & $n(r) = 1 + 2GM/(rc^2)$; symmetric impedance gradient. \\
Spacetime curvature & Dielectric refraction & Geodesic $\equiv$ Fermat path through refractive medium. \\
Stress-energy $T_{\mu\nu}$ & LC energy density $U$ & $U = \tfrac{1}{2}\varepsilon_0 |E|^2 + \tfrac{1}{2}\mu_0 |H|^2$. \\
Schwarzschild radius $r_s$ & Saturation boundary ($S \to 0$) & $Z \to 0, \; \Gamma = -1$: total impedance collapse. \\
Event horizon & Catastrophic impedance boundary & Axiom~4: dielectric yield $\to$ total confinement. \\
Gravitational waves & Transverse inductive shear waves & Low-frequency shear modes in the LC condensate. \\
Frame dragging (Kerr) & Acoustic vortex flow $\vec{v}_\phi$ & Gordon metric: $\omega(r) = 2Mar/(r^2+a^2)^2$. \\
Gravitational lensing & Optical refraction ($n > 1$) & $n(r) = (1+r_s/2r)^3/(1-r_s/2r)$. \\
$G$ (Newton) & TT strain projection & $G = \alpha^2 c^3 \ell_{node}/(28\pi m_e)$ (Axiom~2). \\
``Dark matter'' halo & Regime~I vacuum drag & $v_{flat} = (GM a_0)^{1/4}$; $a_0 = cH_\infty/(2\pi)$. \\
\bottomrule
\end{tabular}
\end{table}


\begin{table}[H]
\centering
\caption{General Relativity $\leftrightarrow$ AVE Translation.  Gravity is macroscopic dielectric refraction; no curved manifold is required.}
\label{tab:trans_gravity}
\small
\renewcommand{\arraystretch}{1.3}
\begin{tabular}{@{}p{3.8cm} p{4.0cm} p{5.8cm}@{}}
\toprule
\textbf{General Relativity} & \textbf{AVE Equivalent} & \textbf{Relationship} \\
\midrule
Metric tensor $g_{\mu\nu}$ & Variable $\varepsilon_{eff}, \mu_{eff}$ & $n(r) = 1 + 2GM/(rc^2)$; symmetric impedance gradient. \\
Spacetime curvature & Dielectric refraction & Geodesic $\equiv$ Fermat path through refractive medium. \\
Stress-energy $T_{\mu\nu}$ & LC energy density $U$ & $U = \tfrac{1}{2}\varepsilon_0 |E|^2 + \tfrac{1}{2}\mu_0 |H|^2$. \\
Schwarzschild radius $r_s$ & Saturation boundary ($S \to 0$) & $Z \to 0, \; \Gamma = -1$: total impedance collapse. \\
Event horizon & Catastrophic impedance boundary & Axiom~4: dielectric yield $\to$ total confinement. \\
Gravitational waves & Transverse inductive shear waves & Low-frequency shear modes in the LC condensate. \\
Frame dragging (Kerr) & Acoustic vortex flow $\vec{v}_\phi$ & Gordon metric: $\omega(r) = 2Mar/(r^2+a^2)^2$. \\
Gravitational lensing & Optical refraction ($n > 1$) & $n(r) = (1+r_s/2r)^3/(1-r_s/2r)$. \\
$G$ (Newton) & TT strain projection & $G = \alpha^2 c^3 \ell_{node}/(28\pi m_e)$ (Axiom~2). \\
``Dark matter'' halo & Regime~I vacuum drag & $v_{flat} = (GM a_0)^{1/4}$; $a_0 = cH_\infty/(2\pi)$. \\
\bottomrule
\end{tabular}
\end{table}


% ═══════════════════════════════════════════════════════════════
% DOMAIN TRANSLATION TABLE: Cosmology ↔ AVE
% Canonical location: manuscript/common/translation_cosmology.tex
% ═══════════════════════════════════════════════════════════════
%
% Usage: % ═══════════════════════════════════════════════════════════════
% DOMAIN TRANSLATION TABLE: Cosmology ↔ AVE
% Canonical location: manuscript/common/translation_cosmology.tex
% ═══════════════════════════════════════════════════════════════
%
% Usage: % ═══════════════════════════════════════════════════════════════
% DOMAIN TRANSLATION TABLE: Cosmology ↔ AVE
% Canonical location: manuscript/common/translation_cosmology.tex
% ═══════════════════════════════════════════════════════════════
%
% Usage: \input{../../common/translation_cosmology.tex}

\begin{table}[H]
\centering
\caption{Cosmology $\leftrightarrow$ AVE Translation. Expansion, dark energy, and dark matter emerge from lattice creep and impedance gradients.}
\label{tab:trans_cosmology}
\small
\renewcommand{\arraystretch}{1.3}
\begin{tabular}{@{}p{3.8cm} p{4.0cm} p{5.8cm}@{}}
\toprule
\textbf{Cosmology} & \textbf{AVE Equivalent} & \textbf{Relationship} \\
\midrule
Hubble expansion $H_0$ & Lattice creep rate $H_\infty$ & $H_\infty = 28\pi m_e^3 cG/(\hbar^2\alpha^2) \approx 69.3$~km/s/Mpc. \\
Dark energy ($\Lambda$) & Latent strain energy & $w_{vac} < -1$: stable phantom from lattice pre-stress. \\
Dark matter & Regime~I vacuum drag & Unsaturated lattice mutual inductance below $a_0$. \\
MOND acceleration $a_0$ & Regime boundary & $a_0 = cH_\infty/(2\pi) \approx 1.07 \times 10^{-10}$~m/s$^2$. \\
Flat rotation curves & Asymptotic visco-kinematic floor & $v_{flat} = (GM_{bar}\,a_0)^{1/4}$ from hoop stress. \\
CMB temperature & Thermal lattice noise $T_{CMB}$ & Residual acoustic background of the LC mesh. \\
Horizon scale $R_H$ & $\alpha^2/(28\pi\alpha_G)$ lattice pitches & $R_H \approx 14.1$~Gly. \\
Inflation & Not required & Lattice isotropy is built in (Axiom~1). \\
Baryon asymmetry & Topological chirality bias & $\mathcal{M}_A$ chiral LC lattice breaks CP natively. \\
\bottomrule
\end{tabular}
\end{table}


\begin{table}[H]
\centering
\caption{Cosmology $\leftrightarrow$ AVE Translation. Expansion, dark energy, and dark matter emerge from lattice creep and impedance gradients.}
\label{tab:trans_cosmology}
\small
\renewcommand{\arraystretch}{1.3}
\begin{tabular}{@{}p{3.8cm} p{4.0cm} p{5.8cm}@{}}
\toprule
\textbf{Cosmology} & \textbf{AVE Equivalent} & \textbf{Relationship} \\
\midrule
Hubble expansion $H_0$ & Lattice creep rate $H_\infty$ & $H_\infty = 28\pi m_e^3 cG/(\hbar^2\alpha^2) \approx 69.3$~km/s/Mpc. \\
Dark energy ($\Lambda$) & Latent strain energy & $w_{vac} < -1$: stable phantom from lattice pre-stress. \\
Dark matter & Regime~I vacuum drag & Unsaturated lattice mutual inductance below $a_0$. \\
MOND acceleration $a_0$ & Regime boundary & $a_0 = cH_\infty/(2\pi) \approx 1.07 \times 10^{-10}$~m/s$^2$. \\
Flat rotation curves & Asymptotic visco-kinematic floor & $v_{flat} = (GM_{bar}\,a_0)^{1/4}$ from hoop stress. \\
CMB temperature & Thermal lattice noise $T_{CMB}$ & Residual acoustic background of the LC mesh. \\
Horizon scale $R_H$ & $\alpha^2/(28\pi\alpha_G)$ lattice pitches & $R_H \approx 14.1$~Gly. \\
Inflation & Not required & Lattice isotropy is built in (Axiom~1). \\
Baryon asymmetry & Topological chirality bias & $\mathcal{M}_A$ chiral LC lattice breaks CP natively. \\
\bottomrule
\end{tabular}
\end{table}


\begin{table}[H]
\centering
\caption{Cosmology $\leftrightarrow$ AVE Translation. Expansion, dark energy, and dark matter emerge from lattice creep and impedance gradients.}
\label{tab:trans_cosmology}
\small
\renewcommand{\arraystretch}{1.3}
\begin{tabular}{@{}p{3.8cm} p{4.0cm} p{5.8cm}@{}}
\toprule
\textbf{Cosmology} & \textbf{AVE Equivalent} & \textbf{Relationship} \\
\midrule
Hubble expansion $H_0$ & Lattice creep rate $H_\infty$ & $H_\infty = 28\pi m_e^3 cG/(\hbar^2\alpha^2) \approx 69.3$~km/s/Mpc. \\
Dark energy ($\Lambda$) & Latent strain energy & $w_{vac} < -1$: stable phantom from lattice pre-stress. \\
Dark matter & Regime~I vacuum drag & Unsaturated lattice mutual inductance below $a_0$. \\
MOND acceleration $a_0$ & Regime boundary & $a_0 = cH_\infty/(2\pi) \approx 1.07 \times 10^{-10}$~m/s$^2$. \\
Flat rotation curves & Asymptotic visco-kinematic floor & $v_{flat} = (GM_{bar}\,a_0)^{1/4}$ from hoop stress. \\
CMB temperature & Thermal lattice noise $T_{CMB}$ & Residual acoustic background of the LC mesh. \\
Horizon scale $R_H$ & $\alpha^2/(28\pi\alpha_G)$ lattice pitches & $R_H \approx 14.1$~Gly. \\
Inflation & Not required & Lattice isotropy is built in (Axiom~1). \\
Baryon asymmetry & Topological chirality bias & $\mathcal{M}_A$ chiral LC lattice breaks CP natively. \\
\bottomrule
\end{tabular}
\end{table}


\begin{summarybox}
\begin{itemize}
    \item General Relativity is formally mapped onto the variable-impedance LC medium, eliminating the need for arbitrary, non-Euclidean geometric manifolds.
    \item Spacetime curvature phenomena, including gravitational lensing and redshift, are completely identical to classical acoustic and optical refractive index fields acting within the dielectric lattice.
    \item Relativistic phenomena such as the D-shaped shadow and relativistic Doppler beaming emerge mechanically from Hamiltonian raymarching through local topological saturation fields.
    \item The K4-TLM native lattice simulator cross-validates gravitational lensing as a native property of a K4 Diamond lattice metric tensor under strain.
\end{itemize}
\end{summarybox}

\begin{exercisebox}
\begin{enumerate}
    \item Using the isotropic-coordinate redshift factor $z = W/U$, calculate the suppression of emission for an accretion disk annulus located uniformly at $r = 1.1 \, r_h$. 
    \item Demonstrate mathematically how the Schwarzschild coordinates map exactly to the saturated Diamond lattice impedance metric via Op14.
\end{enumerate}
\end{exercisebox}

